\section{Preliminary Investigation: Token-Efficient Projection Routing}
\label{sec:projection}

The 100-case ablation shows that retrieval policy, not only model capability, determines whether the orchestrator receives useful evidence. To optimize this bottleneck, we implemented a Context Projection Model trained as a persona-conditioned cross-encoder. The projector operates after Master Context Sphere assembly and before persona prompting, scoring candidate graph nodes separately for the Product Manager, Worker, and Reviewer.

\subsection{Training Methodology}

The projection model was trained on 7{,}299 training samples with an 888-row held-out validation split. Training used the \texttt{cross-encoder/ms-marco-MiniLM-L-6-v2} backbone on Apple Silicon MPS. To address the severe class imbalance in implementation-level routing, we applied persona-stratified oversampling, increasing Worker-positive training rows from 109 to 279 while preserving the grouped validation split. The model was optimized with asymmetric binary cross-entropy using persona-specific positive weights (Worker=18, Reviewer=10, PM=8) and a negative weight of 1.0.

We selected the final checkpoint using a Worker Margin criterion, defined as the validation mean score for Worker-positive rows minus the validation mean score for Worker-negative rows. This early-stopping rule selected the Epoch 1 checkpoint, before later optimization could overfit easier PM/Reviewer structural labels and degrade the fragile Worker routing signal. The released artifact repository includes the training scripts, saved report, thresholds, validation scores, and model checkpoint metadata.

\subsection{Asymmetric Cognitive Routing}

The projection task is asymmetric across personas. Product Manager and Reviewer prompts benefit from high-precision structural skeletons: paths, interfaces, tests, imports, and risk-bearing adjacency. Worker prompts require higher-recall code access because implementation often depends on local function bodies and nearby edit anchors. Initial training exposed gradient starvation for implementation-level routing, where Worker-positive evidence was underrepresented relative to easier structural negatives.

\subsection{Safety Floors Against Context Starvation}

Aggressive thresholding reduced token load but created a failure mode we call \textit{context starvation}: a persona could receive too few nodes to maintain a coherent local view. In the no-floor projection smoke test, the Reviewer often received zero selected nodes, and the pass rate dropped to 6/10 on a cohort that full Context Sphere solved completely. We therefore added a deterministic Top-\(K\) safety floor with \(min\_k=2\), ensuring that every persona receives at least two context anchors even when calibrated model confidence is low.

\subsection{Projection Smoke Results}

Table~\ref{tab:projection_results} summarizes the projection smoke test. The cohort consists of 10 cases selected from Context Sphere successes in the 100-case ablation, so the baseline pass rate is intentionally 10/10. With the Top-\(K\) safety floor, projection preserved 9/10 passes while reducing input tokens by 71.5\% and estimated inference cost by 58.4\%.

\begin{table}[t]
\centering
\caption{Projection-mode routing on a 10-case Context Sphere success cohort.}
\label{tab:projection_results}
\resizebox{\columnwidth}{!}{%
\begin{tabular}{lccc}
\toprule
\textbf{Metric} & \textbf{Full Context Sphere} & \textbf{Projection (\(min\_k=2\))} & \textbf{Change} \\
\midrule
Local pass rate & 10 / 10 & 9 / 10 & -10.0\% \\
Total tokens & 2{,}090{,}215 & 643{,}512 & \textbf{-69.2\%} \\
Input tokens & 2{,}078{,}983 & 592{,}935 & \textbf{-71.5\%} \\
Context characters & 8{,}483{,}908 & 1{,}979{,}745 & \textbf{-76.7\%} \\
Inference cost & \$0.5737 & \$0.2386 & \textbf{-58.4\%} \\
\bottomrule
\end{tabular}%
}
\end{table}

This result reframes projection as a resource-control layer rather than a replacement for the Context Sphere. The Master Sphere remains responsible for broad structural recall; projection converts that master topology into persona-specific Spheres of Influence that constrain token load while preserving enough evidence for downstream repair.
